% ==========================================================================
% Paper 2 — Section 6: From Routing to Binding. Representation Decoupling
% (TAN-II Sprint 4.2-C, frozen).
% ==========================================================================
\section{From Routing to Binding: Representation Decoupling}
\label{sec:binding}

Routing changes \emph{which} candidate is selected; binding additionally
requires that the selected candidate's \emph{content} be transmitted, and
that the transmission follow the (key, value) association rather than key
identity, position, or amplitude. The minimal degree of freedom that
separates the two is operator O3 of §2.4: each slot now carries a pair
$(k_i,v_i)$, keys drive the attention, values are transmitted, and the
query slots carry $v=0$ and are never candidates. This section reports the
frozen Sprint 4.2-C audits and states the structural obstruction---the
softmax mixing barrier---that forces the next rung.

\subsection{The binding probe}
\label{sec:probe6}

The frozen canonical history is $[A=(1,5),\,0,\,B=(2,9),\,0]$ (keys
inherited from §5, values $V_A=5$, $V_B=9$), the counterfactual queries are
the heritage pair $(3.0,4.0)$, and the readout is the attended value
$C=\sum_i\alpha_i v_i$, with the events-only normalization
$C_{\mathrm{ev}}=(\alpha_Av_A+\alpha_Bv_B)/(\alpha_A+\alpha_B)$ reported
alongside. Three quantities are pre-registered: the binding errors
$E_{\mathrm{bind},A}=|\hat C_1-V_A|$, $E_{\mathrm{bind},B}=|\hat
C_2-V_B|$; the counterfactual transfer
$T=C_{\mathrm{ev}}(Q_A)-C_{\mathrm{ev}}(Q_B)$; and the softness distance
$D=|C_{\mathrm{ev}}-v_{\mathrm{winner}}|$. The value-swap control
$(1,9),(2,5)$ is the identifiability instrument: a system that transmits
the \emph{selected} value must satisfy $T_{\mathrm{swap}}=-T$ exactly,
while a system that merely reads fixed magnitudes cannot.

\subsection{Soft binding, exactly pairing-sensitive}
\label{sec:softresults}

The frozen deterministic results are:

\begin{itemize}
\item \emph{Winner-aligned transfer.} $C_{\mathrm{ev}}(3)=6.7427$ and
$C_{\mathrm{ev}}(4)=7.1556$: the readout moves toward the selected value at
both queries ($6.7427$ closer to $V_A=5$ than to $9$; $7.1556$ closer to
$V_B=9$ than to $5$), with $T=-0.4129$ carrying the sign of
$v_{\mathrm{winner}}(Q_A)-v_{\mathrm{winner}}(Q_B)$.
\item \emph{Exact value-swap antisymmetry.} $T_{\mathrm{swap}}=+0.4129$,
with $T_{\mathrm{swap}}=-T$ verified to $10^{-12}$: the output follows the
(key,value) association, not the key identity and not the amplitude.
\item \emph{The softness is real.} $D=1.7427$ at $Q=3$ and $1.8444$ at
$Q=4$: the delivered value is a convex mixture of both candidates, not the
bound symbol.
\item \emph{The baselines separate routing from sharpening.} The scalar
model (winner always $B$) shows only a sharpening drift
$T_{\mathrm{M0}}=-0.0844$, five times smaller than the routing model's
transfer; the sign control M1 ($s=\pm1$) binds \emph{harder} on this
history ($D_{\mathrm{M1}}=0.1064$, energies spread widely) yet remains the
degenerate order-locked reversal of §5; and the unnormalized control
collapses onto the probe slot and transmits nothing (Amendment AM2).
\item \emph{Ensemble.} Over randomized keys (values fixed by identity,
positions randomized) the transfer magnitude averages $0.4004$
(quadrature benchmark; Monte Carlo $0.3977$, $0.3967$, $0.4021$), with
$E[|T|\mid\text{flip}]=0.4655$ versus $E[|T|\mid\text{no flip}]=0.3483$
and the sign-matched alignment on flips exactly $1$---the latter a
theorem-level check (the event-normalized share is a logistic function of
the logit difference), reported as such rather than as a discovery.
\end{itemize}

\subsection{The softmax mixing barrier}
\label{sec:mixingbarrier}

The softness distance is not a numerical shortfall; it is a structural
consequence of the readout.

\begin{proposition}[Softmax mixing barrier]
\label{prop:mixing}
Let $V_A \neq V_B$ be distinct candidate values. At any finite sharpening $\gamma<\infty$,
the attention weights satisfy $\alpha_i>0$ for every candidate with finite logits.
Under the event-normalized readout $C_{\mathrm{ev}} = (\alpha_A V_A + \alpha_B V_B)/(\alpha_A + \alpha_B)$
(with $\alpha_A+\alpha_B>0$), $C_{\mathrm{ev}}$ lies in the strict relative interior of the segment
connecting $V_A$ and $V_B$: $C_{\mathrm{ev}} = \lambda V_A + (1-\lambda)V_B$ with
$\lambda = \alpha_A/(\alpha_A+\alpha_B) \in (0, 1)$. Consequently, the softness distance satisfies
\begin{equation}
D_{\mathrm{ev}} = \min(|C_{\mathrm{ev}}-V_A|,\, |C_{\mathrm{ev}}-V_B|) = \min(\lambda, 1-\lambda)|V_A - V_B| > 0.
\end{equation}
Exact value delivery ($D_{\mathrm{ev}}=0$) strictly requires a discrete selection step that sets the
losing weight to zero: continuous softmax weighting over distinct candidates, at any finite temperature,
cannot deliver an isolated bound symbol.
\end{proposition}

\begin{remark}[Cancellation artifacts in unnormalized full-window readouts]
\label{rem:cancellation}
When unnormalized full-window readouts $C_{\mathrm{full}} = \sum_{i=1}^W \alpha_i v_i$ are evaluated over
sequences containing silent background slots ($v_{\mathrm{silent}}=0$), the active candidate weights sum to
$\sum_{i\in\mathrm{events}}\alpha_i < 1$. This dilution allows $C_{\mathrm{full}}$ to pass through $V_A$
at an intermediate temperature ($\gamma \approx 1.198$, documented in Amendment~AM5), yielding a spurious
zero $|C_{\mathrm{full}} - V_A| = 0$. This zero is a cancellation artifact of background dilution, not hard
symbolic binding; under the event-normalized readout $C_{\mathrm{ev}}$, softness remains $D_{\mathrm{ev}} \approx 1.5 > 0$.
\end{remark}

Proposition~\ref{prop:mixing} converts the measured
$D\approx1.7$--$1.8$ from a ``performance'' observation into a boundary
statement of the same rank as Theorem~\ref{thm:order}: \emph{soft attention
$\neq$ exact value selection}. It also explains the one honest inversion
found in the frozen transmission audits: through the excitatory branch
alone the drive transfer keeps the value-aligned direction
($-0.2586$), but through the full E-I circuit the inhibitory branch
re-reads the same value channel and subtracts it, inverting the net
transfer ($+0.2367$; frozen Amendment AM4 corrects the numeric prediction,
the direction statement is unchanged). E-I organization therefore does not
implement binding; it \emph{reshapes the value-transfer pathway after
query-conditioned routing}---an observation carried into §7, where the
selection step is added.

\subsection{Status of the rung}
\label{sec:status6}

The frozen verdicts read: binding at the readout level is \emph{soft but
genuine} (winner-aligned, pairing-sensitive to $10^{-12}$, baselines
separated), and exact delivery is \emph{not} achieved (the mixing barrier,
$D>0$). Routing and binding are now experimentally decoupled: §5 supplies
the reordering, this section supplies the value channel, and the missing
ingredient---discrete selection---is the single degree of freedom added in
§7.
